\chapter{What this manual cannot state, and where the documents are wrong}\label{ch:unproven}

A manual that does not say what it has not tested is a manual a developer trusts in the wrong place.
This chapter is the list of things this manual asserts only as ``the code says'', not as ``we saw it
happen''.

\section{Never run here}

\begin{itemize}
  \item \textbf{Real hardware.} Every Device test in the suite talks to something in the same
    process: the \texttt{LoopbackPeer} and \texttt{UdpLoopbackPeer} in \texttt{tests/\_fixtures.py},
    a socket on an ephemeral port, a pseudo-terminal for \texttt{SERIAL}, and pyvisa-py behind a
    \texttt{TCPIP0::127.0.0.1::\dots::SOCKET} resource for \texttt{VISA}. No GPIB instrument, no real
    serial adapter, no NI-VISA stack has ever been opened by this suite. The driver contract is
    proven; the driver against an instrument is not.
  \item \textbf{Real Windows sessions.} The same tests run on \texttt{windows-latest} in CI, and
    the notes record those CI runs being green. Nothing in this manual was executed on Windows by
    us. The \texttt{icacls} permission paths, the Inno installer and the frozen Windows bundle are
    asserted by reading the scripts and by \texttt{tests/test\_release\_ci\_27c.py}, which checks the
    workflow's shape, never by an install.
  \item \textbf{The Debian container smokes, on this machine.}
    \texttt{ATE\_DOCKER\_TESTS=1 pytest -m build} and \texttt{packaging/linux/docker\_smoke.py} were
    attempted while writing this manual and did not complete: the container ran out of scratch space.
    \texttt{tests/test\_release\_ci\_28c.py} asserts the gate is written as documented and the smoke
    scripts' logic; the containers themselves are CI's job. Nothing in this manual claims those
    smokes passed here.
  \item \textbf{A real Chromium PDF.} \texttt{conftest.py} replaces the PDF helper with a stub for
    every test except the four that must prove a PDF has text. The sandbox decision, the flags and
    the failure messages are tested; the rendered page is not, except in those four.
  \item \textbf{Both Qt bindings.} The PySide6/PyQt6 differences in chapter~\ref{ch:pyside} come
    from the 32a switch and its notes. They were not measured side by side on this machine; where a
    claim rests on one binding only, ADR-0001 says so.
\end{itemize}

\section{The 32b stall: fixed, not explained}

Packet 32b is the clearest example of a fix that outran its diagnosis. One Windows CI run left a
Test Run stalled for tens of seconds. The packet added the abort-boundary handling, the
\texttt{describe\_stuck\_runs} report, and the store's write timing, and the tests prove each of
them behaves. The cause --- that the disk was the thing that was slow --- is not established. The
diagnostic exists so that the next stall can be read: long or in-flight writes point at the disk,
normal writes with the worker's stack elsewhere point at something nobody has found. A manual that
stated the cause as fact would be the same kind of error as a test that widens a timeout to make a
failure disappear.

\section{Where the documents and the code disagree}

These were found while writing the manuals. In every case the code was followed and the document was
left as it is, so that the difference is visible rather than smoothed over:

\begin{itemize}
  \item \textbf{Spec rule 20 says \texttt{ate db migrate} upgrades ``after a timestamped backup''.}
    It makes no backup. \texttt{cli.\_cmd\_db} calls \texttt{RunStore}, which creates the missing
    tables and sets \texttt{PRAGMA user\_version}, and prints
    \texttt{created or upgraded database at \dots}. The admin manual tells an administrator to back
    up first, and says so in the same breath as the command.
  \item \textbf{Spec rule 08 says \texttt{log\_line} rows are written ``about once a second''.}
    They are inserted per step, not on a timer. The review notes recorded this in packet 08 and it
    has never been corrected in the spec.
  \item \textbf{The release workflow has three jobs, not five.} The 32a notes and the review notes
    both say ``all 5 jobs''; \texttt{release.yml} defines \texttt{windows-installer},
    \texttt{linux-build} and \texttt{linux-smoke}, and \texttt{linux-smoke} runs a matrix of
    \texttt{deb}, \texttt{tar} and \texttt{wheel} --- five job runs from three job definitions. The
    count is right; the word is not.
  \item \textbf{\texttt{\_location()} finds a step's file and line by searching the schedule text}
    for \texttt{:setAction(}, not from the Lua call stack. For a multi-line Action, a helper module
    or a Sub-schedule the reported location can be wrong. The manual describes what is reported, not
    what the notes imply it should be.
  \item \textbf{The prototype layout in \texttt{run-screen-prototype.html} is not the screen.} The
  staged Run screen is a vertical stack; the review notes say so, and the developer manual describes
  the widgets the code creates.
\end{itemize}

\section{What changes the moment you read this}

Two numbers in this manual are snapshots: the suite's size and the elapsed time. Both were taken on
one machine on one day, and both move with every packet. The names, the signals, the tables and the
drivers are not snapshots --- the suite derives them from the code and fails if a manual stops naming
them. If \texttt{1571} is no longer the number, the tests in \texttt{tests/test\_manuals.py} do not
care; if a signal is gone from \texttt{RunEngine} and the manual still lists it, nothing notices, so
the direction of the risk matters: the checks catch a manual falling behind, not a manual ahead of
the code.

\cando{say what a claim in this manual rests on, and what would falsify it.}
